\subsection{Tổ chức chương trình}
\subsubsection{Kiến trúc tổng thể của chương trình}
Hệ thống được thiết kế và xây dựng dựa trên mô hình Lập trình theo mô-đun (module). Thay vì tập trung toàn bộ mã nguồn vào một tệp duy nhất, chương trình được phân rã thành các mô-đun độc lập dựa trên \textit{nguyên lý Phân tách trách nhiệm}.

Cách tiếp cận này đảm bảo hệ thống đạt được độ gắn kết cao bên trong từng mô-đun và độ phụ thuộc thấp giữa các mô-đun với nhau, tạo tiền đề thuận lợi cho việc kiểm thử, gỡ lỗi và mở rộng tính năng trong tương lai.
\subsubsection{Cấu trúc tệp tin và phân rã chức năng}

Mã nguồn của đồ án được tổ chức thành 6 thành phần cốt lõi, bao gồm 1 tệp tiêu đề cơ sở, 4 cặp module logic (\texttt{.h} và \texttt{.cpp}) và một tệp điều khiển trung tâm. Chi tiết chức năng của từng tệp được quy định như sau:
\begin{enumerate}[label*=(\alph*)]
	\item \texttt{DataStructures.h}: Định nghĩa cấu trúc nút LCRS, bản ghi và biến toàn cục, đảm bảo nhất quán kiểu dữ liệu giữa các mô-đun.
	\item \texttt{TreeManager.h/.cpp}: Thao tác trực tiếp trên RAM; khởi tạo biến toàn cục; quản lý vòng đời nút (cấp phát, giải phóng, truy xuất).
	\item \texttt{SearchEngine.h/.cpp}: Xử lý truy vấn, kiểm soát trùng lặp; duyệt DFS để tìm kiếm; cung cấp giao diện chọn thực thể.
	\item \texttt{Relationship.h/.cpp}: Giải quyết bài toán đồ thị gia phả; tìm tổ tiên chung (LCA); đổi tọa độ thế hệ sang danh xưng và ngược lại.
	\item \texttt{FileHandler.h/.cpp}: Quản lý I/O với ổ cứng; tuần tự hóa cây LCRS thành mảng, ghi/đọc file nhị phân; hỗ trợ nhập từ file văn bản, xử lý lỗi dòng.
	\item \texttt{main.cpp}: Khởi chạy, tạo dữ liệu mẫu, điều hướng sự kiện và xử lý lỗi nhập.
\end{enumerate}

\begin{figure}[h!]
	\begin{tikzpicture}[
		box/.style={draw, rounded corners=6pt, minimum width=2.8cm, minimum height=1cm, align=center, font=\small},
		 >={Stealth},
		mainbox/.style={box, fill=yellow!30, draw=yellow!50!black, very thick},
		bluebox/.style={box, fill=blue!20, draw=blue!50!black, very thick},
		redbox/.style={box, fill=red!20, draw=red!60!black, very thick},
		arrow/.style={->, very thick},
		darrow/.style={<->, very thick},
		dottedarrow/.style={->, dotted, very thick},
		]
		
		% Main
		\node[mainbox] (main) at (0,0) {main.cpp};
		
		% Middle layer (FIXED positions)
		\node[bluebox] (search) at (-6,-2.5) {SearchEngine};
		\node[bluebox] (file)   at (-2,-2.5) {FileHandler};
		\node[bluebox] (rel)    at (2,-2.5)  {Relationship};
		\node[bluebox] (tree)   at (6,-2.5)  {TreeManager};
		
		% Bottom
		\node[redbox] (data) at (0,-5) {DataStructures};
		
		% Arrows
		\draw[arrow] (main) -- (search);
		\draw[arrow] (main) -- (file);
		\draw[arrow] (main) -- (rel);
		\draw[arrow] (main) -- (tree);
		
		\draw[darrow] (search) -- (file);
		\draw[darrow] (file) -- (rel);
		\draw[darrow] (rel) -- (tree);
		
		\draw[dottedarrow] (search) -- (data);
		\draw[dottedarrow] (file) -- (data);
		\draw[dottedarrow] (rel) -- (data);
		\draw[dottedarrow] (tree) -- (data);
		\draw[dottedarrow] (main) -- (data);
		
		\node[draw, rectangle, right=2cm of data, align=right, scale=0.8] (legend) {
			\begin{tikzpicture}[scale=0.9]
				\draw[darrow] (0,0) -- (1,0) node[right]{Hỗ trợ};
				\draw[arrow] (0,-0.5) -- (1,-0.5) node[right]{Điều khiển};
				\draw[dottedarrow] (0,-1) -- (1,-1) node[right]{Dựa trên nền tảng của};
			\end{tikzpicture}
		};
		
	\end{tikzpicture}
	\caption{Sơ đồ mô tả cấu trúc chương trình và các tệp tin}
\end{figure}

\begin{sidewaysfigure}
	\begin{tikzpicture}[
		node distance=1cm,
		block/.style={rectangle, draw=blue!70, fill=blue!5, thick, minimum width=3cm, minimum height=0.8cm, align=center, font=\bfseries\small, rounded corners},
		subblock/.style={rectangle, draw=gray!60, fill=gray!5, minimum width=2.5cm, minimum height=0.7cm, align=center, font=\footnotesize},
		mainblock/.style={rectangle, draw=red!70, fill=red!10, thick, minimum width=4cm, minimum height=1cm, align=center, font=\bfseries},
		managerblock/.style={rectangle, draw=orange!70, fill=orange!5, thick, minimum width=3.5cm, minimum height=0.8cm, align=center, font=\bfseries\small},
		arrow/.style={-stealth, thick, draw=gray!80}
		]
		
		% --- TẦNG 1: ĐIỀU KHIỂN ---
		\node[mainblock] (main) {main.cpp \\ \scriptsize (Vòng lặp điều hướng)};
		
		% --- TẦNG 2: MODULE LOGIC (Dàn hàng ngang rộng) ---
		\node[block, below=of main] (rel) {Relationship.cpp \\ \scriptsize (Phân tích quan hệ)};
		\node[block, left=5cm of rel] (search) {SearchEngine.cpp \\ \scriptsize (Tìm kiếm thực thể)};
		\node[block, right=5cm of rel] (file) {FileHandler.cpp \\ \scriptsize (Xử lý tệp tin)};
		
		% --- TẦNG 3: CÁC HÀM CHI TIẾT ---
		% Nhánh Search
		\node[subblock, below left=0.8cm and -0.5cm of search] (dfs) {SearchByName() \\ \scriptsize (Thuật toán DFS)};
		\node[subblock, below right=0.8cm and -0.5cm of search] (hash) {IDLookup() \\ \scriptsize (Dùng Bảng băm)};
		
		% Nhánh Relationship
		\node[subblock, below=1cm of rel] (lca) {DetermineRel() \\ \scriptsize (Thuật toán LCA)};
		\node[subblock, left=0.3cm of lca] (path) {GetPathToRoot()};
		\node[subblock, right=0.3cm of lca] (anc) {FindAncestor()};
		\node[subblock, below=0.6cm of anc] (coll) {CollectCands() \\ \scriptsize (Tìm theo tầng)};
		
		% Nhánh File
		\node[subblock, below left=0.8cm and -0.4cm of file] (save) {SaveFile() \\ \scriptsize (Serialization)};
		\node[subblock, below right=0.8cm and -0.4cm of file] (load) {LoadFile() \\ \scriptsize (Reconstruct)};
		
		% --- TẦNG 4: TREE MANAGER (Dùng chung) ---
		\node[managerblock, below=4.5cm of rel] (manager) {TreeManager.cpp \\ \scriptsize (Dịch vụ bộ nhớ RAM)};
		\node[subblock, left=0.5cm of manager] (add) {AddChild()};
		\node[subblock, right=0.5cm of manager] (free) {FreeTree()};
		
		% --- KẾT NỐI MŨI TÊN ---
		% Main gọi 3 module chính
		\draw[arrow] (main.south) -- (rel.north);
		\draw[arrow] (main.south) -- ++(0,-0.6) -| (search.north);
		\draw[arrow] (main.south) -- ++(0,-0.6) -| (file.north);
		
		% Nội bộ Search
		\draw[arrow] (search.south) -- ++(0,-0.4) -| (dfs.north);
		\draw[arrow] (search.south) -- ++(0,-0.4) -| (hash.north);
		
		% Nội bộ Relationship
		\draw[arrow] (rel.south) -- (lca.north);
		\draw[arrow] (rel.south) -- ++(0,-0.5) -| (path.north);
		\draw[arrow] (rel.south) -- ++(0,-0.5) -| (anc.north);
		\draw[arrow] (anc.south) -- (coll.north);
		
		% Nội bộ File
		\draw[arrow] (file.south) -- ++(0,-0.4) -| (save.north);
		\draw[arrow] (file.south) -- ++(0,-0.4) -| (load.north);
		
		% Tất cả đổ về TreeManager khi cần thao tác nút
		\coordinate (mid) at (0,-6.8);
		\draw[arrow, dashed] (dfs.south) |- (mid) -| (manager.north);
		\draw[arrow, dashed] (lca.south) -- (manager.north);
		\draw[arrow, dashed] (load.south) |- (mid) -| (manager.north);
		
		\draw[arrow] (manager.west) -- (add.east);
		\draw[arrow] (manager.east) -- (free.west);
		
	\end{tikzpicture}
	\caption{Sơ đồ cấu trúc các hàm chính của chương trình}
\end{sidewaysfigure}

Toàn bộ mã nguồn của chương trình quản lý gia phả có tại Phụ lục đính kèm theo báo cáo này.

\subsection{Ngôn ngữ cài đặt}

\subsubsection{Lý do sử dụng ngôn ngữ C++}
Dự án lựa chọn ngôn ngữ C++ làm ngôn ngữ lập trình chính để triển khai hệ thống bởi những lý do như sau:
\begin{itemize}
	\item Mô hình LCRS yêu cầu cấp phát bộ nhớ liên tục, C++ dùng new/delete để kiểm soát vòng đời của từng nút (Person) một cách chính xác, nhờ đó tối ưu hiệu suất.
	\item Dùng \texttt{std::string} thay vì mảng char giúp xử lý tên tiếng Việt an toàn, tránh tràn bộ đệm. \texttt{std::vector} cũng được dùng để lưu kết quả tìm kiếm, giúp mã nguồn ngắn gọn, rõ ràng.
	\item Nhờ struct/class, C++ mô hình hóa “Con người” tự nhiên với đầy đủ thuộc tính và quan hệ, đồng thời dễ dàng tách chương trình thành các tệp \texttt{.h} và \texttt{.cpp}.
\end{itemize}

\subsubsection{Môi trường phát triển và công cụ hỗ trợ}
\begin{itemize}
	\item Trình biên dịch: Sử dụng GCC (\texttt{g++}) tiêu chuẩn C++11 trở lên để hỗ trợ các tính năng hiện đại và đảm bảo tốc độ thực thi tối ưu.
	\item Hệ thống được phát triển trên các IDE phổ biến như \textit{Embarcadero Dev-C++} và \textit{Visual Studio}, tận dụng khả năng quản lý \textit{project} để liên kết các tệp tin mô-đun.
	\item Định dạng tệp lưu trữ: tệp nhị phân (.dat) dùng để lưu trữ trạng thái cây dưới dạng các struct tĩnh, đảm bảo tốc độ đọc/ghi nhanh và bảo mật dữ liệu cơ bản; tệp văn bản (\texttt{.txt}) Dùng để nhập liệu hàng loạt theo định dạng phân tách bằng dấu phẩy (CSV), giúp người dùng dễ dàng soạn thảo dữ liệu từ bên ngoài.
\end{itemize}

\subsection{Kết quả}